% TOP_PROBLEM: {"id": 30003894, "title": "Arithmetic Nature of Delocalized Eta Invariants", "category_id": 16, "category": {"id": 16, "name": "physics", "display_name": "Mathematical Physics", "description": "Problems at the intersection of mathematics and physics.", "slug": "physics", "order_index": 16, "created_at": "2026-07-31T15:26:25.671Z"}, "status": "open", "statusQualification": "The group, conjugacy class, operator and convergence hypotheses are essential to define the invariant. The source does not ask for algebraicity of divergent or undefined expressions, and the general value question is broader than the highlighted algebraicity question."}
% FIRST_POSTED: 2026-10-11
% ENTRY_KIND: statement_only
% SUBMISSION: {"kind": "statement_only", "category": "top_problems", "problem_id": "30003894", "proposal_type": "established", "source": "Non-commutative Geometry, Index Theory and Mathematical Physics (2018), Oberwolfach Reports, DOI 10.4171/owr/2018/32, pp. 1911--1981. Primary report accessed 2026-10-11; relevant PDF pages 12. https://doi.org/10.4171/owr/2018/32; Xie--Yu, IMRN (catalogue-listed context; not independently consulted in this preparation). https://doi.org/10.1093/imrn/rnz170", "status_evidence": "Canonical source status open retained. Supplied literature triage dated 2026-08-21: The question whether delocalized eta invariants are always algebraic remains open in general; algebraicity is proved conditionally from rational Baum--Connes and related hypotheses. Primary-report access checked 2026-10-11; this historical report and inherited triage do not certify current resolution of every component.", "status_checked_on": "2026-10-11", "duplicate_check": "Compared canonical ID, title and normalized target with the complete local UnsolvedMath registry; normalized exact-text matches: none; title matches: none. Related or overlapping canonical records: none identified in this batch. Primary source context and the existing OWR statement audit used for scope and duplicate review; an exhaustive mathematical-equivalence/literature search was not performed.", "attribution": "Original proposer not named in the supplied canonical record; attribution remains to the identified Oberwolfach contribution/report, not its editors as invented proposers. Explanatory statement preparation and source/semantic audit assisted by OpenAI GPT-6 Codex; no human mathematical authorship or independent proof verification is claimed."}
% SOURCE_STATEMENT: "What values can delocalized eta invariants of Dirac operators on universal covers take in general; in particular, are they always algebraic numbers?"
% Preparation and final self-audit: CONTRIBUTING.md prompts C and D; source/audit artifacts retained in the local physics_statements_2026_10_11 batch.
% Imported from: physics_statements.tex; UnsolvedMath ID 30003894
% !TeX program = lualatex
\documentclass[11pt]{article}
\usepackage[margin=25mm]{geometry}
\usepackage{fontspec}
\setmainfont{Latin Modern Roman}
\usepackage{amsmath,amssymb,mathtools}
\usepackage{unicode-math}
\setmathfont{Latin Modern Math}
\usepackage{xurl,hyperref}
\hypersetup{colorlinks=true,urlcolor=blue}
\setcounter{secnumdepth}{0}
\setlength{\parindent}{0pt}
\setlength{\parskip}{5pt}
\setlength{\emergencystretch}{3em}
\newcommand{\sref}[2]{\hyperlink{source:#1:#2}{[S#2]}}
\newcommand{\cataloguescope}[1]{\paragraph{Recorded target.}#1}
\newcommand{\R}{\mathbb R}
\newcommand{\C}{\mathbb C}
\newcommand{\N}{\mathbb N}
\newcommand{\Z}{\mathbb Z}
\newcommand{\Q}{\mathbb Q}
\newcommand{\F}{\mathbb F}
\newcommand{\eps}{\varepsilon}
\begin{document}
% UnsolvedMath ID 30003894; source catalogue code OWR-16407-005
\setcounter{section}{30003893}
\section{Arithmetic Nature of Delocalized Eta Invariants}\label{30003894}
{\small UnsolvedMath ID 30003894 · Mathematical Physics}
\subsection{Definitions and mathematical statement}

\paragraph{Shared notation.}$\mathbb N=\{1,2,\ldots\}$; the standard number fields have their usual meanings. All additional conventions are stated locally.


\paragraph{Definitions and conventions.}
For an odd-dimensional Dirac operator $D$, an eta invariant measures spectral asymmetry through a regularized trace integral. On a universal cover with deck group $\Gamma$, a delocalized trace sums the integral kernel over a specified nontrivial conjugacy class of $\Gamma$. Replacing the usual trace by this trace gives a delocalized eta invariant when the summation and regularization converge. An algebraic number is a root of a nonzero polynomial with rational coefficients.

\paragraph{Target retained from the canonical source catalogue.}
\paragraph{Catalogue wording retained for traceability.}
What values can delocalized eta invariants of Dirac operators on universal covers take in general; in particular, are they always algebraic numbers?

\paragraph{Source scope and qualification.}
The group, conjugacy class, operator and convergence hypotheses are essential to define the invariant. The source does not ask for algebraicity of divergent or undefined expressions, and the general value question is broader than the highlighted algebraicity question.
\subsection{Short English statement}
Which numbers can delocalized spectral-asymmetry invariants take, and must they always be algebraic?
\subsection{Sources}
\begin{itemize}
\item[\hypertarget{source:30003894:1}{S1}] ULAM.AI and the UnsolvedMath Contributors, UnsolvedMath, record 30003894 (OWR-16407-005). Catalogue curation under CC BY 4.0; definitions and English exposition added locally.
\url{https://huggingface.co/datasets/ulamai/UnsolvedMath}.
\item[\hypertarget{source:30003894:2}{S2}] Non-commutative Geometry, Index Theory and Mathematical Physics (2018), Oberwolfach Reports, DOI 10.4171/owr/2018/32, pp. 1911--1981. Primary report accessed 2026-10-11; relevant PDF pages 12.
\url{https://doi.org/10.4171/owr/2018/32}.
\item[\hypertarget{source:30003894:3}{S3}] Xie--Yu, IMRN (catalogue-listed context; not independently consulted in this preparation).
\url{https://doi.org/10.1093/imrn/rnz170}.
\end{itemize}
\label{end:30003894}
\end{document}
